\section{Discussion, Limitations and Adversarial Cues}
\label{sec:discussion}

\paragraph{What is established and what is not.} The results of \cref{sec:properties} concern the \emph{decision}. That every trace of the emitted Tamarin theory uses only values usable in the corresponding cognitive state is not proved; the host's certification and the invariant of RQ5 check given models, but an invariant of emitted instrumentation misses a dropped lapse marker, a fabricated recovery event, or a dropped branch. A translation theorem for the supported fragment, or a validator with a soundness proof~\cite{pnueli1998translation}, is the main open obligation; completeness with forgotten derivations (\cref{prop:bounded}) is the second.

\paragraph{Where the current framework stops.} In every compiled theory, lapses are defeaters and no rule supports $\neg\Via$; the only conflicts are between a default and its exceptions, resolved by fixed priorities. This is why \cref{thm:char} has a closed form, why the ambiguity regime is irrelevant, and why usability is monotone in cues: a cue can only help. But attackers control many surfaces on which cues appear, such as a spoofed login page or an altered invoice, and memory is distorted by misleading information~\cite{loftus2005planting} and competing retrievals~\cite{anderson1994remembering}.

\paragraph{A preliminary model of adversarial cues.} We extend the cue model with an \emph{interference relation} $\Int\subseteq\Cue_A\times K$, a set $\Pres_A$ of adversarial cues present, and a strict partial order $\succ$ on all cues, read as salience or trust. The compilation adds, for every present $(c_a,h)\in\Int$ and every derivation $\pi$ with $h\in\Hyp_\pi$, a fact $\AdvCue(c_a,h)$, facts $\mathit{ContextCue}(c,h)$ for $c\in\mathit{Pres}$ with $(c,h)\in\mathit{Assoc}$, and an \emph{interference rule}
\[
\iota_{\pi,h,c_a}:\ \AdvCue(c_a,h)\defeas\neg\Via(\pi,t),
\]
emits reinstaters $\rho_{\pi,h,c}$ also for interfered hypotheses, with $\rho_{\pi,h,c}>\iota_{\pi,h,c_a}$ if $c\succ c_a$ and $\iota_{\pi,h,c_a}>\rho_{\pi,h,c}$ if $c_a\succ c$ ($\iota$ and $d_\pi$ are unordered), and lets interfered hypotheses generate substitutes like blocked ones. Since $\iota$ is a defeasible rule, not a defeater, it creates a two-sided conflict with $d_\pi$. To express the decision inside the logic, we add $s_t:\Usable(t)\defeas\Send(t)$ for every candidate and $p_{t'}:\Usable(m_2)\defeas\neg\Send(t')$ with $p_{t'}>s_{t'}$ for every substitute $t'$; the decision is the $t$ with $+\mathit{Send}(t)$, or Skip if there is none. Without adversarial cues this layer is conservative under $\partial$: the $\Send$ tags give exactly the decision of \cref{alg:decision} (proof in \cref{app:proofs}; checked on 3{,}000 random theories).

\begin{table}[t]
\centering\scriptsize
\caption{Preliminary adversarial-cue model on the running example (computed with the tool's reasoner, $\partial$, and the ambiguity-propagating prototype, $\delta$). The banner interferes with $\pwa$; the hint is associated with $\pwa$.}
\label{tab:adversarial}
\setlength{\tabcolsep}{3pt}
\begin{tabular}{@{}llccc@{}}
\toprule
 & Situation & $\Usable(\mthree)$, $\partial$ & decision, $\partial$ & decision, $\delta$ \\
\midrule
A0 & no lapse, no cue & $+$ & \preserve & \preserve \\
A1 & no lapse, banner & $-$ & \replace{} $m_3'$ & \skipm \\
A2 & no lapse, banner and hint, hint $\succ$ banner & $+$ & \preserve & \preserve \\
A3 & $\pwa$ forgotten, banner and hint, unordered & $-$ & \replace{} $m_3'$ & \skipm \\
A4 & $\pwa$ forgotten, banner and hint, banner $\succ$ hint & $-$ & \replace{} $m_3'$ & \replace{} $m_3'$ \\
A5 & $\pwa$ forgotten, banner and hint, hint $\succ$ banner & $+$ & \preserve & \preserve \\
\bottomrule
\end{tabular}
\end{table}

\Cref{tab:adversarial} shows what changes. First, usability is no longer monotone in the cues present (A0 versus A1): an interfering cue alone makes a remembered password unusable, and under $\partial$ the human sends the other bank's password, which enables the attack of Sect. 2 without any memory lapse. Second, the conflict between remembering and interference is resolved by priorities: a trusted cue that outranks the adversarial one restores the intended send (A2, A5). In general, if every present adversarial cue on a hypothesis is outranked by a present genuine cue associated with it, the tags of all $\Via$ and $\Usable$ atoms are those of the theory without adversarial cues (\cref{app:proofs}). Third, when the conflict is unresolved (A1, A3), the choice of semantics decides the outcome. Under ambiguity blocking, the doubtful $\Usable(\mthree)$ is refuted, the preference for the intended message is discarded, and the human sends the alternative; under ambiguity propagation, $\Usable(\mthree)$ remains supported, the preference still blocks the alternative, and nothing is sent. The two semantics thus encode two human policies, acting on what remains versus withholding when in doubt, and the framework turns this choice into an explicit, analysable parameter. Here the priority $\omega>d_\pi$, inert for $\partial$ (\cref{prop:priorities}), matters: under propagation it decides whether a forgotten value still creates doubt. In A4 the propagating semantics replaces with the pairs $\omega>d_\pi$ and skips without them.

These results are computed by the existing reasoners on extended theories built by the research harness; the compiler does not produce them, and nothing here reaches Tamarin. A full treatment needs cues taken from attacker-controlled inputs, a mechanism for recalling a specific wrong value, which interference alone does not provide, and verification goals such as ``no redirection under trusted-cue precedence'', on ceremonies such as payments where a provider-generated confirmation competes with details printed on an attacker-controlled document.